\documentclass[12pt, a4paper]{article}
\usepackage[utf8]{inputenc}
\usepackage[lithuanian]{babel}
\usepackage[T1]{fontenc}
\usepackage{amsmath, amssymb}
\usepackage{geometry}
\usepackage{enumitem}
\usepackage{tasks}
\settasks{label-width=2em, item-indent=3em}

\geometry{top=2cm, bottom=2cm, left=2cm, right=2cm}

\title{\textbf{Namų darbas: algebrinis valsas}}
\author{}
\date{2026-10-03}

\begin{document}

\maketitle

\section*{Namų darbo atlikimo tvarka}

\begin{itemize}[leftmargin=*]
    \item Užduotis atlikite atskiruose A4 formato lapuose ir atneškite trečiadienį. Darbai bus surinkti.
    \item Sprendimus rašykite įskaitomai ir tvarkingai. Neįskaitomas, netvarkingas arba pribraukytas darbas nebus priimtas – jį reikės perrašyti ir tą pačią dieną atnešti.
    \item Atsakymus pažymėkite spalva arba apibrėžkite stačiakampiu.
    \item Jei užstrigsite, klauskite per „Messenger“ arba „Teams“ – padėsiu. Pirmadienį per pamokas taip pat galėsite spręsti šias užduotis ir klausti.
    \item Nepalikite darbo paskutinei dienai.
    \item Su Vijeto teorema žiūrėkite pavyzdžius ir spręskite užduotis pagal juos. 
\end{itemize}

\section{Sąlygos}

\begin{enumerate}[label=\textbf{\arabic* užduotis.}, leftmargin=*]
    \item Apskaičiuokite reiškinio reikšmę:
    \begin{tasks}[label={(\arabic*)}, label-offset=0.9em](1)
        \task $\dfrac{(3\sqrt{2}-2\sqrt{5})(\sqrt{7}-\sqrt{3})(3\sqrt{2}+2\sqrt{5})(\sqrt{7}+\sqrt{3})}{-\text{DBD}(104,136)}$
        \task $\dfrac{(\sqrt{28}-\sqrt{12})^{666}\cdot(2\sqrt{7}+2\sqrt{3})^{666}\cdot 0{,}0625^{666}}{\dfrac{1}{\pi^{1+2+\ldots+20}}}\cdot\pi^{-208}$
        \task $\left(\frac{1}{1+\frac{1}{1+\frac{1}{2}}}\right)^{-1}:\left(3+\frac{1}{1+\frac{1}{8}}\right)\cdot 2{,}\overline{3}$
    \end{tasks}

    \item Supaprastinkite:
    \begin{tasks}[label={(\arabic*)}, label-offset=0.9em](1)
        \task $\dfrac{16x^4+54x}{8x^3-18x} \div \dfrac{4x^4-6x^3+9x^2}{4x^4-9x^2}$
        \task $\dfrac{27x^6-8x^3}{9x^5-4x^3} \div \dfrac{18x^4+12x^3+8x^2}{18x^4-8x^2}$
        \task $\dfrac{x^3-64}{x^2+8x+16} \cdot \dfrac{x^2-16}{x^2+4x+16} \div \dfrac{x-4}{x+4}$
        \task $\dfrac{a^5 b+a^2 b^4}{a^4 b-2 a^3 b^2+a^2 b^3} \div \dfrac{a^4 b^2+a b^5}{a^3 b^2-a b^4} \cdot \dfrac{2 a^4-2 a b^3}{2 a^3+2 a^2 b+2 a b^2}$
    \end{tasks}

    \item Raskite didžiausią daugianario
    $$P(x)=(2x-1)^3-4(x+1)(2-x)(3-2x)-7\left(1\frac{2}{7}\right)x^2$$
    reikšmę.

    \item Neskaičiuodami lygties sprendinių, raskite simetrinių reiškinių reikšmes (žr. Vjeto teoremos pavyzdžius 6–9):
    \begin{tasks}[label={(\arabic*)}, label-offset=0.9em](1)
        \task Lygties $x^2-4x+3=0$ sprendiniai yra $x_1$ ir $x_2$. Raskite $x_1^2+x_2^2$.
        \task Lygties $-x^2+3x-2=0$ sprendiniai yra $x_1$ ir $x_2$. Raskite $x_1^3+x_2^3$.
        \task Lygties $-x^2-5x+3=0$ sprendiniai yra $x_1$ ir $x_2$. Raskite $\dfrac{1}{x_1}+\dfrac{1}{x_2}$.
        \task Lygties $x^2+2x-1=0$ sprendiniai yra $x_1$ ir $x_2$. Raskite $x_1^2x_2+x_1x_2^2$.
        \task Lygties $-x^2-2x+3=0$ sprendiniai yra $x_1$ ir $x_2$. Raskite $(x_1-x_2)^2$.
        \task Lygties $x^2+4x+2=0$ sprendiniai yra $x_1$ ir $x_2$. Raskite $x_1^2x_2^2$.
        \task Lygties $-x^2-6x+1=0$ sprendiniai yra $x_1$ ir $x_2$. Raskite $\dfrac{x_1}{x_2}+\dfrac{x_2}{x_1}$.
        \task Lygties $x^2+x-5=0$ sprendiniai yra $x_1$ ir $x_2$. Raskite $\dfrac{1}{x_1^2}+\dfrac{1}{x_2^2}$.
        \task Lygties $-x^2-3x+1=0$ sprendiniai yra $x_1$ ir $x_2$. Raskite $x_1^3x_2+x_1x_2^3$.
        \task Lygties $x^2+5x+4=0$ sprendiniai yra $x_1$ ir $x_2$. Raskite $\dfrac{1}{x_1^2}+\dfrac{1}{x_2^2}$.
    \end{tasks}

    \item Neskaičiuodami lygties sprendinių, raskite nurodytų reiškinių reikšmes (pirmiesiems penkiems punktams žiūrėkite pavyzdžius 6–9, kitiems penkiems – pavyzdžius 12–15):
    \begin{tasks}[label={(\arabic*)}, label-offset=0.9em](1)
        \task Lygties $x^2-4x+1=0$ sprendiniai yra $x_1$ ir $x_2$. Raskite $3x_1^2+2x_2^2-10x_1-6x_2+10$.
        \task Lygties $x^2+5x+2=0$ sprendiniai yra $x_1$ ir $x_2$. Raskite $x_1^2-2x_2^2+8x_1-7x_2-6$.
        \task Lygties $x^2-3x-2=0$ sprendiniai yra $x_1$ ir $x_2$. Raskite $4x_1^2+x_2^2-13x_1-4x_2$.
        \task Lygties $x^2+x-5=0$ sprendiniai yra $x_1$ ir $x_2$. Raskite $-2x_1^2+3x_2^2+2x_1+7x_2-4$.
        \task Lygties $x^2-7x+3=0$ sprendiniai yra $x_1$ ir $x_2$. Raskite $2x_1^2+2x_2^2-13x_1-13x_2+10$.
    \end{tasks}
    \begin{tasks}[label={(\arabic*)}, label-offset=0.9em, resume](1)
        \task Lygties $x^2-2x-1=0$ sprendiniai yra $x_1$ ir $x_2$. Raskite $x_1^3+2x_2^3-5x_2^2+5x_2+2$.
        \task Lygties $x^2+3x+1=0$ sprendiniai yra $x_1$ ir $x_2$. Raskite $2x_1^3+x_2^3+5x_1^2+6x_2^2+x_1+12x_2-2$.
        \task Lygties $x^2-5x+2=0$ sprendiniai yra $x_1$ ir $x_2$. Raskite $x_1^3-x_2^3-7x_1^2+6x_2^2+15x_1-4x_2-1$.
        \task Lygties $x^2-x-4=0$ sprendiniai yra $x_1$ ir $x_2$. Raskite $2x_1^3+2x_2^3-x_1^2-5x_2^2-11x_1-7x_2+15$.
        \task Lygties $x^2+4x-3=0$ sprendiniai yra $x_1$ ir $x_2$. Raskite $x_1^3+x_2^3+4x_1^2+2x_2^2-2x_1-10x_2+6$.
    \end{tasks}
\end{enumerate}

\newpage
\section{Atsakymai}

\begin{enumerate}[label=\textbf{\arabic* užduotis.}, leftmargin=*]
    \item
    \begin{tasks}[label={(\arabic*)}, label-offset=0.9em](1)
        \task $1$
        \task $\pi^2$
        \task $1$
    \end{tasks}
    \item
    \begin{tasks}[label={(\arabic*)}, label-offset=0.9em](1)
        \task $2x+3$
        \task $3x-2$
        \task $x-4$
        \task $a+b$
    \end{tasks}
    \item $0$ (kai $x=5$).
    \item
    \begin{tasks}[label={(\arabic*)}, label-offset=0.9em](1)
        \task $10$
        \task $9$
        \task $\dfrac{5}{3}$
        \task $2$
        \task $16$
        \task $4$
        \task $-38$
        \task $\dfrac{11}{25}$
        \task $-11$
        \task $\dfrac{17}{16}$
    \end{tasks}
    \item
    \begin{tasks}[label={(\arabic*)}, label-offset=0.9em](1)
        \task $13$
        \task $-19$
        \task $7$
        \task $-3$
        \task $5$
        \task $13$
        \task $-10$
        \task $2$
        \task $5$
        \task $-4$
    \end{tasks}
\end{enumerate}

\end{document}
